%========= SURVEY MACROS ===================================================
\usepackage[table]{xcolor}
\usepackage{tikz}
\usepackage{pgf-pie}
\usepackage{pgfplots}
\usepackage{multicol} 
\usepackage{capt-of} 
\pgfplotsset{compat=1.18}
\usepackage{xfp} 

%========= STUDENT BAR CHART MACRO ===================================================
\newcommand{\studentsurveybar}[6]{
    \def\totalresponses{\fpeval{#1 + #2 + #3 + #4 + #5}}
    \def\meanweight{\fpeval{round((1*#1 + 2*#2 + 3*#3 + 4*#4 + 5*#5) / \totalresponses, 2)}}

    \begin{center}
    \begin{minipage}{\linewidth} %
        \vspace{-1em} % Reduces gap above chart
        \begin{tikzpicture}
            \begin{axis}[
                ybar, bar width=0.6cm, ymin=0,
                enlarge y limits={upper, value=0.2},
                width=\linewidth, height=4.5cm, % Slightly shorter to fit more on page
                axis y line*=left, ymajorgrids=true, ytick distance=20,
                symbolic x coords={1,2,3,4,5}, xtick=data, enlarge x limits=0.2,
                nodes near coords,
                every node near coord/.append style={
                    font=\tiny, color=white, anchor=north, yshift=-2pt
                }
            ]
            \addplot[fill=brown!60!black, draw=none] coordinates {
                (1,#1) [#1 (\fpeval{round(#1/\totalresponses*100, 1)}\%)]
                (2,#2) [#2 (\fpeval{round(#2/\totalresponses*100, 1)}\%)]
                (3,#3) [#3 (\fpeval{round(#3/\totalresponses*100, 1)}\%)]
                (4,#4) [#4 (\fpeval{round(#4/\totalresponses*100, 1)}\%)]
                (5,#5) [#5 (\fpeval{round(#5/\totalresponses*100, 1)}\%)]
            };
            \end{axis}
        \end{tikzpicture}
        \vspace{-0.5em}
        \captionof{figure}{#6 (Mean: \meanweight)}
        \vspace{0.5em} % Small buffer before next chart
    \end{minipage}
    \end{center}
}

%========= BAR CHART MACRO ===================================================
\newcommand{\surveybar}[6]{
    \def\totalresponses{\fpeval{#1 + #2 + #3 + #4 + #5}}
    \def\meanweight{\fpeval{round((1*#1 + 2*#2 + 3*#3 + 4*#4 + 5*#5) / \totalresponses, 2)}}

    \begin{center}
    \begin{minipage}{\linewidth} %
        \vspace{-1em} % Reduces gap above chart
        \begin{tikzpicture}
            \begin{axis}[
                ybar, bar width=0.6cm, ymin=0,
                enlarge y limits={upper, value=0.2},
                width=\linewidth, height=4.5cm, % Slightly shorter to fit more on page
                axis y line*=left, ymajorgrids=true, ytick distance=1,
                symbolic x coords={1,2,3,4,5}, xtick=data, enlarge x limits=0.2,
                nodes near coords,
                every node near coord/.append style={
                    font=\tiny, color=white, anchor=north, yshift=-2pt
                }
            ]
            \addplot[fill=brown!60!black, draw=none] coordinates {
                (1,#1) [#1 (\fpeval{round(#1/\totalresponses*100, 1)}\%)]
                (2,#2) [#2 (\fpeval{round(#2/\totalresponses*100, 1)}\%)]
                (3,#3) [#3 (\fpeval{round(#3/\totalresponses*100, 1)}\%)]
                (4,#4) [#4 (\fpeval{round(#4/\totalresponses*100, 1)}\%)]
                (5,#5) [#5 (\fpeval{round(#5/\totalresponses*100, 1)}\%)]
            };
            \end{axis}
        \end{tikzpicture}
        \vspace{-0.5em}
        \captionof{figure}{#6 (Mean: \meanweight)}
        \vspace{0.5em} % Small buffer before next chart
    \end{minipage}
    \end{center}
}

%========= PIE CHART MACRO ===================================================
\newcommand{\surveypie}[2]{
    \begin{center}
    \begin{minipage}{\linewidth} %
        \begin{tikzpicture}
            \pie[sum = auto, text = legend, radius = 1.5]{#1}
        \end{tikzpicture}
        \captionof{figure}{#2}
    \end{minipage}
    \end{center}
}

%========= CONTINUOUS SECTION MACRO =========================================
\newcommand{\surveysection}[1]{
    \vspace{2.5em}           % Single top margin for the section
    \hrule
    \nopagebreak             % Prevents the line from being separated from the title
    \vspace{0.6em}           % Space between the line and the text
    \noindent{\large\textbf{#1}}
    \par\nopagebreak         % Ensures the title stays with the chart below
    \vspace{0.8em}           % Space before the first chart starts
}
